\paragraph{端点三态压缩原理：\texorpdfstring{$\mathcal S_{\star}=\{0,2,-2\}$}{S*={0,2,-2}} 与模 \texorpdfstring{$4$}{4} 三通道}
为把命题 \ref{prop:half-angle-z4-residue} 的 \(\ZZ/4\) 残差进一步压缩为“有限状态”模板，引入端点三态集合。

\begin{definition}[端点三态集合与端点态函数]\label{def:endpoint-tristate-set}
在端点 \(w=-1\)（等价于 \(u=-1\)，且 \(s=0\)）处，记
\[
\boxed{\ \mathcal S_{\star}:=\{0,2,-2\}\ }.
\]
并称映射
\[
d\ \longmapsto\ C_d(0)\in\mathcal S_{\star}\qquad(d\ge 1)
\]
为幂伸缩指数 \(d\) 的\textbf{端点态函数}。
\end{definition}

\begin{corollary}[端点三态压缩原理（\texorpdfstring{$\bmod 4$}{mod 4} 三通道退化）]\label{cor:endpoint-tristate-compression}
对一切整数 \(d\ge 1\)，端点取值满足
\begin{equation}\label{eq:endpoint-tristate-values}
\boxed{\;
C_d(0)=\mathrm{i}^{d}+\mathrm{i}^{-d}=2\cos(d\pi/2)\in\mathcal S_{\star},\qquad
C_d(0)=
\begin{cases}
0, & d\ \text{为奇数},\\
2, & d\equiv 0\ (\mathrm{mod}\ 4),\\
-2,& d\equiv 2\ (\mathrm{mod}\ 4).
\end{cases}\;}
\end{equation}
因此，任意由幂替换半群 \(\{u\mapsto u^d\}_{d\ge 1}\) 诱导的完成化动力学 \(s\mapsto C_d(s)\) 在端点 \(s=0\) 的 \(0\)-喷射上必然退化为一个模 \(4\) 的三态系统。
\end{corollary}
\begin{proof}
式 \eqref{eq:endpoint-tristate-values} 即命题 \ref{prop:half-angle-z4-residue} 的端点恒等式 \eqref{eq:chebyshev-endpoint-values} 在三类剩余 \(d\bmod 4\) 下的展开写法。证毕。
\end{proof}

\paragraph{Witt--Chebyshev 端点喷射分解：primitive 的三态模板}
令 \(p_n(u)\in\ZZ[u]\) 为 primitive 多项式，并在完成化坐标 \(u=r^2,\ s=r+r^{-1}\) 下定义
\(\widehat p_n(s):=r^{-n}p_n(r^2)\in\ZZ[s]\)；
同时对任意 \(d\ge 1\) 记 \(C_d(s):=r^d+r^{-d}\in\ZZ[s]\)。
则有严格倍角律 \(r^{-dn}p_n(r^{2d})=\widehat p_n(C_d(s))\)（见推论 \ref{cor:primitive-completion-hatp}）。
据此引入端点三态值
\begin{equation}\label{eq:primitive-endpoint-tristate-values}
\boxed{\;
A_n:=\widehat p_n(0),\qquad
B_n^{+}:=\widehat p_n(2),\qquad
B_n^{-}:=\widehat p_n(-2)\; }.
\end{equation}

\begin{proposition}[primitive 的端点三态退化：\texorpdfstring{$\widehat p_n(C_d(0))$}{hat p_n(C_d(0))} 的模 \texorpdfstring{$4$}{4} 通道分裂]\label{prop:primitive-endpoint-tristate-degeneration}
对任意 \(n\ge 1\) 与 \(d\ge 1\)，在端点 \(s=0\) 有严格退化
\begin{equation}\label{eq:primitive-endpoint-tristate-degeneration}
\boxed{\;
\widehat p_n\!\bigl(C_d(0)\bigr)=
\begin{cases}
A_n, & d \text{ 奇},\\
B_n^{+}, & d\equiv 0\ (\mathrm{mod}\ 4),\\
B_n^{-}, & d\equiv 2\ (\mathrm{mod}\ 4).
\end{cases}\;}
\end{equation}
并且端点两侧 \(\pm 2\) 与 \(u=1\) 的关系给出
\begin{equation}\label{eq:primitive-endpoint-Bpm-identification}
\boxed{\ B_n^{+}=p_n(1),\qquad B_n^{-}=(-1)^n p_n(1)\ }.
\end{equation}
\end{proposition}
\begin{proof}
由推论 \ref{cor:endpoint-tristate-compression}，端点态 \(C_d(0)\) 仅取三值 \(\{0,2,-2\}\)，且其取值由 \(d\bmod 4\) 决定；代入 \(\widehat p_n\) 即得 \eqref{eq:primitive-endpoint-tristate-degeneration}。
\par
另一方面，由定义 \(\widehat p_n(s)=r^{-n}p_n(r^2)\)，当 \(s=2\) 时取 \(r=1\) 得 \(\widehat p_n(2)=p_n(1)\)。
当 \(s=-2\) 时取 \(r=-1\) 得 \(\widehat p_n(-2)=(-1)^{-n}p_n(1)=(-1)^n p_n(1)\)。
从而得到 \eqref{eq:primitive-endpoint-Bpm-identification}。证毕。
\end{proof}

\paragraph{端点 WITT 分解模板：三通道权重核与对数主尺度}
对在包含区间 \([-2,2]\) 的某开邻域上解析的函数 \(g\)，定义截断的完成化幂链算子
\begin{equation}\label{eq:completed-witt-pow-operator-trunc}
\boxed{\;
(W_{\mathrm{pow}}^{\le M}g)(s):=\sum_{d=1}^{M}\frac{1}{d}\,g\!\bigl(C_d(s)\bigr)\qquad(M\ge 1)\; }.
\end{equation}
则端点 \(s=0\) 的 \(0\)-喷射满足严格的模类通道分解：

\begin{proposition}[端点三通道分解：截断幂链在 \texorpdfstring{$s=0$}{s=0} 的模 \texorpdfstring{$4$}{4} 分组恒等式]\label{prop:endpoint-mod4-channel-decomposition}
令
\[
H_M^{(1)}:=\sum_{\substack{1\le d\le M\\ d\ \mathrm{odd}}}\frac1d,\qquad
H_M^{(0)}:=\sum_{\substack{1\le d\le M\\ d\equiv 0\ (\mathrm{mod}\ 4)}}\frac1d,\qquad
H_M^{(2)}:=\sum_{\substack{1\le d\le M\\ d\equiv 2\ (\mathrm{mod}\ 4)}}\frac1d.
\]
则对任意解析 \(g\) 有恒等式
\begin{equation}\label{eq:endpoint-mod4-channel-decomposition}
\boxed{\;
(W_{\mathrm{pow}}^{\le M}g)(0)
=
H_M^{(1)}\,g(0)
+H_M^{(0)}\,g(2)
+H_M^{(2)}\,g(-2)\; }.
\end{equation}
因此端点处的幂链叠加在 \(0\)-喷射层面退化为三个互不耦合的模类通道，权重完全由 \(\sum_{d\equiv a\ (\mathrm{mod}\ 4)}\frac1d\) 型核决定。
\end{proposition}
\begin{proof}
由定义 \eqref{eq:completed-witt-pow-operator-trunc} 与端点三态压缩 \eqref{eq:endpoint-tristate-values}，
对每一项 \(g(C_d(0))\) 按 \(d\bmod 4\) 分类并合并即可得到 \eqref{eq:endpoint-mod4-channel-decomposition}。证毕。
\end{proof}

\begin{remark}[端点对数主尺度与相位精度解耦：\texorpdfstring{$u=-1$}{u=-1} 下偶幂通道的计数控制]\label{rem:endpoint-logscale-parity-decouple}
在 twist 层 \(u=-1\) 上，偶幂满足 \(u^{2k}=1\)，故任何 Witt 型伸缩叠加
\(
\sum_{m\ge 1}\frac1m f(u^m)
\)
都包含一个由偶幂子链 \(\sum_{k\ge 1}\frac{1}{2k}f(1)\) 诱导的对数主尺度。
对 primitive 输入 \(f=p_n\) 而言，该主尺度的系数为
\(
f(1)=p_n(1)=B_n^{+}
\)
（见 \eqref{eq:primitive-endpoint-Bpm-identification}），其与端点差分轴
\(
A_n=\widehat p_n(0)
\)
所编码的半角分支信息在定义上分离。
特别地，当 \(n\) 为奇数时 \(\widehat p_n\) 为奇多项式（推论 \ref{cor:primitive-completion-hatp}），从而 \(A_n=\widehat p_n(0)=0\)；而当 \(n\ge 3\) 为奇数时还满足 \(p_n(-1)=0\)（推论 \ref{cor:primitive-odd-vanish-u-minus1}），端点奇偶差分轴在奇长 primitive 上完全消失。
\end{remark}

\paragraph{Frobenius--Adams 整数层相位证书：\texorpdfstring{$\Delta_{p,k}$}{Delta} 在端点三态上的分层审计}
定义 \ref{def:witt-frobenius-defect} 给出 Frobenius 缺陷多项式
\(
\Delta_{p,k}(u):=a_{p^k}(u)-a_{p^{k-1}}(u^p)\in\ZZ[u]
\)，
并由命题 \ref{prop:witt-frobenius-defect-divisibility} 得到整除深度
\(
\Delta_{p,k}(u)\in p^k\ZZ[u]
\)。
另一方面，系数环版本的 Big-Witt/Möbius 反演在带 Adams 操作族 \(\psi^d\) 的交换环 \(R\) 上满足
\(
n\,p_n=\sum_{d\mid n}\mu(d)\,\psi^d(a_{n/d})\in R
\)
（命题 \ref{prop:witt-necklace-dwork-R}），其中在多参数多项式环情形 \(\psi^d\) 由逐变量幂替换实现。
将上述两条与端点三态压缩并置，可定义一个可直接审计实现链路的整数证书族。

\begin{definition}[端点三态上的 Frobenius 整数证书]\label{def:endpoint-frobenius-integer-certificate}
对任意素数 \(p\)、整数 \(k\ge 1\) 与端点态 \(\sigma\in\mathcal S_{\star}\)，取 \(r_\sigma\in\{\,\mathrm{i},1,-1\,\}\) 使 \(r_\sigma+r_\sigma^{-1}=\sigma\)，并令 \(u_\sigma:=r_\sigma^{2}\in\{-1,1\}\)。
定义
\begin{equation}\label{eq:endpoint-frobenius-integer-certificate}
\boxed{\ \mathfrak I_{p,k}(\sigma):=\frac{\Delta_{p,k}(u_{\sigma})}{p^{k}}\in\ZZ\ }.
\end{equation}
\end{definition}

\begin{proposition}[整数层先验护栏：端点证书的可审计性]\label{prop:endpoint-frobenius-integer-certificate-auditable}
对任意 \((p,k)\) 与 \(\sigma\in\mathcal S_{\star}\)，证书 \(\mathfrak I_{p,k}(\sigma)\) 为良定义的整数。
若在任一实现中出现 \(\mathfrak I_{p,k}(\sigma)\notin\ZZ\)（等价于 \(\Delta_{p,k}(u_\sigma)\notin p^k\ZZ\)），则必违反 Frobenius 缺陷整除深度 \(\Delta_{p,k}(u)\in p^k\ZZ[u]\)，从而可在整数层先行定位错误，再进入单位圆谱半径等连续谱层讨论。
\end{proposition}
\begin{proof}
由命题 \ref{prop:witt-frobenius-defect-divisibility}，\(\Delta_{p,k}(u)\) 的各系数都被 \(p^k\) 整除；因此对任意 \(u_\sigma\in\ZZ\) 有 \(\Delta_{p,k}(u_\sigma)\in p^k\ZZ\)，从而 \(\mathfrak I_{p,k}(\sigma)\in\ZZ\)。
若出现反例，则与“系数级整除”矛盾。证毕。
\end{proof}

\paragraph{整数层残差与谱层偏移的正交性：Frobenius 缺陷作为算术精度轴}
由 \(\Delta_{p,k}(u)\in p^k\ZZ[u]\) 可在整数层定义规范化残差多项式
\begin{equation}\label{eq:frobenius-normalized-residual-poly}
\boxed{\ E_{p,k}(u):=\frac{\Delta_{p,k}(u)}{p^{k}}\in\ZZ[u]\ }.
\end{equation}
于是 \(\{E_{p,k}\}_{p,k}\) 构成一个完全离散的“算术精度轴”：其审计只依赖整数系数与整除深度，而与单位圆谱半径、端点精度势等连续谱层指标无关。
在端点三态 \(\sigma\in\mathcal S_{\star}\) 上，上述残差与定义 \ref{def:endpoint-frobenius-integer-certificate} 的证书严格一致：
\[
\boxed{\ \mathfrak I_{p,k}(\sigma)=E_{p,k}(u_\sigma)\in\ZZ\ }.
\]
从而“先整数层、后谱层”的分层审计原则可被表述为：先以 \eqref{eq:frobenius-normalized-residual-poly}（或其端点取值）锁定实现的一致性护栏，
再进入端点精度势 \(y=-\log|\mathfrak{w}|\)（式 \eqref{eq:endpoint-busemann-potential}）与 unitary slice 上的连续谱偏移等分析层讨论。

